\documentclass[12pt,a4paper]{ctexart}

\usepackage[a4paper,margin=2.2cm]{geometry}
\usepackage{graphicx}
\usepackage{float}
\usepackage{booktabs}
\usepackage{array}
\usepackage{hyperref}
\usepackage{xcolor}
\usepackage{enumitem}
\usepackage{caption}
\usepackage{subcaption}
\usepackage{fancyhdr}
\usepackage{lastpage}
\usepackage{listings}
\usepackage{tikz}
\usetikzlibrary{arrows.meta,positioning,shapes.geometric}

\hypersetup{
  colorlinks=true,
  linkcolor=blue,
  urlcolor=blue,
  citecolor=blue
}

\pagestyle{fancy}
\fancyhf{}
\setlength{\headheight}{15pt}
\fancyfoot[C]{第 \thepage 页，共 \pageref{LastPage} 页}

\lstset{
  basicstyle=\ttfamily\small,
  frame=single,
  breaklines=true,
  numbers=left,
  numberstyle=\tiny,
  keywordstyle=\color{blue},
  commentstyle=\color{teal},
  showstringspaces=false
}

% 以下信息建议学生按自己的实际情况直接修改。
\newcommand{\course}{嵌入式系统设计}
\newcommand{\homeworkname}{作业1 GPIO}
\newcommand{\studentname}{孙凯}
\newcommand{\studentid}{3124001830}
\newcommand{\classname}{微电子 4 班}
\newcommand{\submitdate}{\today}

\fancyhead[L]{\course}
\fancyhead[R]{\homeworkname}

\newcommand{\placeholderfigure}[2]{
  \begin{figure}[H]
    \centering
    \IfFileExists{#1}{
      \includegraphics[width=0.95\linewidth]{#1}
    }{
      \fbox{\parbox[c][5cm][c]{0.9\linewidth}{\centering
      请将图片文件放入 \texttt{#1}\\
      然后重新编译文档}}
    }
    \caption{#2}
  \end{figure}
}

\begin{document}

\pagenumbering{arabic}
\setcounter{tocdepth}{1}

\begin{center}
  {\zihao{3}\bfseries \homeworkname \par}
  \vspace{0.8em}
  \begin{tabular}{ccc}
    \studentname & \studentid & \classname
  \end{tabular}
  \\
  \vspace{0.4em}
  \submitdate
\end{center}

\vspace{0.5em}
\tableofcontents
\vspace{0.8em}

\section{作业目标与任务要求}

本节简要说明本次作业要完成的功能、验收标准。

    1.  实现功能：按键按下时 LED 亮，松开时 LED 灭。
    
    2.  硬件资源和软件工具：STM32F103C8T6 开发板,ST-Link,USB 数据线、杜邦线、按键、LED。
    
    3.  结果判定标准：按键按下时LED不会出现闪烁，能够稳定发光。
    

\section{实验环境}

\subsection{硬件环境}

\begin{itemize}
  \item 开发板：STM32F103C8T6 开发板
  \item 调试器：ST-Link
  \item 连接器件：USB 数据线、杜邦线、按键、LED 等
\end{itemize}

\subsection{软件环境}

\begin{itemize}
  \item STM32CubeIDE
\end{itemize}

\section{实验原理与方案说明}

本节应结合具体作业内容，说明所用外设或中间件的原理。

\begin{enumerate}[label=\arabic*.]
  \item 硬件连接关系：LED 负极 → PA7，LED 正极 → 3.3V ，按键一端 →                                             PB12，按键另一端 → GND 。
  \item 外设配置原理

        PA7    模式：GPIO\_Output  输出类型：Push-Pull（推挽输出） 
                                上拉 / 下拉：无   

                                推挽输出可以稳定输出高 / 低电平，控制LED 亮灭

        PB12    模式：GPIO\_Input     上拉 / 下拉：Pull-up（内部上拉）

                                开启内部上拉后，按键未按下时引脚保持稳定高电平；按下                                  按键后引脚被拉低 
  \item 程序执行逻辑：系统初始化 → 循环检测按键状态 → 根据按键状态控制                                           LED 
\end{enumerate}

\section{STM32CubeMX 配置过程}

关键配置步骤：

\begin{enumerate}[label=\arabic*.]
  \item 芯片型号选择：STM32F103C8T6 
  \item RCC 时钟配置：使用内部 HSI，8 MHz 
  \item SYS 调试接口配置：Serial Wire 
  \item GPIO外设配置：如图

\end{enumerate}

 
\placeholderfigure{figures/cubemx-config.png}{CubeMX 关键配置截图}

\section{程序设计与关键代码说明}

\subsection{主程序结构}

\begin{lstlisting}[language=C,caption={主循环关键代码示例}]
while (1)
{
  GPIO_PinState key_state = HAL_GPIO_ReadPin(GPIOB,         GPIO_PIN_12);
  HAL_GPIO_WritePin(GPIOA, GPIO_PIN_7, key_state);

}
\end{lstlisting}
功能说明：第一个函数读取按键状态，第二个函数根据按键状态控制PA7电平高低，进而控制灯的亮灭。

\placeholderfigure{figures/code-main.png}{关键代码截图}

\section{程序流程图}



\begin{figure}[H]
  \centering
  \begin{tikzpicture}[
    node distance=10mm and 18mm,
    every node/.style={font=\small},
    startstop/.style={rectangle, rounded corners, minimum width=28mm, minimum height=9mm, draw=black, fill=gray!15},
    process/.style={rectangle, minimum width=34mm, minimum height=9mm, draw=black, fill=blue!8},
    decision/.style={diamond, aspect=2, draw=black, fill=orange!12, inner sep=1pt},
    line/.style={-Latex, thick}
  ]
    \node[startstop] (start) {开始};
    \node[process, below=of start] (init) {系统初始化};
    \node[process, below=of init] (task) {执行主功能};
    \node[decision, below=of task, yshift=-2mm] (judge) {按键是否按下};
    \node[process, below left=of judge] (branch1) {LED亮};
    \node[process, below right=of judge] (branch2) {LED灭};
    \node[startstop, below=18mm of judge] (endloop) {循环};

    \draw[line] (start) -- (init);
    \draw[line] (init) -- (task);
    \draw[line] (task) -- (judge);
    \draw[line] (judge) -- node[left] {是} (branch1);
    \draw[line] (judge) -- node[right] {否} (branch2);
    \draw[line] (branch1) |- (endloop);
    \draw[line] (branch2) |- (endloop);
  \end{tikzpicture}
  \caption{程序流程图示例}
\end{figure}

\section{实验结果与现象分析}



实验现象描述：按下按键时LED亮，不松手时LED持续亮，松开按键LED灭。与预期结果一致。


\placeholderfigure{figures/result-photo.jpg}{实验现象照片}

\placeholderfigure{figures/result-serial.png}{运行结果截图或上位机交互截图}

\section{问题分析与调试记录}

暂无

\section{思考题或课后作业回答}

1.  输入端口PB12，输入按键状态，输出端口PA7，输出高低电平控制LED的亮灭，GND->按键->PB12，PA7->LED->3.3V

2.PB12上拉输入原因：未按按键时，PB12为高电平，按下按键时，PB12为低电平

3.PA7推挽输出：可以输出高低电平，驱动能力强，低速输出低能耗，可靠性强

4.结果与预期一致

\section{总结}

掌握了通过cubemx配置rcc，sys，gpio；能够正确读取gpio口的电平状态；能够使LED准确稳定发光；后续可以尝试控制多个LED亮灭。


\section{附录}

可在附录中补充：

\begin{itemize}
  \item 完整关键代码。
  \item 更多截图或照片。
  \item 参考资料。
  \item 自己完成的扩展任务说明。
\end{itemize}

\end{document}